% Generated draft: requires mathematical and source review.
\documentclass[11pt]{article}
\usepackage[margin=25mm]{geometry}
\usepackage{fontspec}
\setmainfont{Latin Modern Roman}
\usepackage{amsmath,amssymb,mathtools,unicode-math}
\setmathfont{Latin Modern Math}
\usepackage{enumitem,needspace,xurl,hyperref}
\setlength{\parindent}{0pt}
\setlength{\parskip}{5pt}
\setlength{\emergencystretch}{3em}
\newcommand{\N}{\mathbb{N}}
\newcommand{\Z}{\mathbb{Z}}
\newcommand{\Q}{\mathbb{Q}}
\newcommand{\R}{\mathbb{R}}
\newcommand{\C}{\mathbb{C}}
\newcommand{\F}{\mathbb{F}}
\newcommand{\A}{\mathbb{A}}
\newcommand{\PP}{\mathbb{P}}
\newcommand{\E}{\mathbb{E}}
\newcommand{\eps}{\varepsilon}
\newcommand{\dd}{\,\mathrm{d}}
\begin{document}
\section*{UnsolvedMath ID: 701604}
\textbf{Source problem code:} catalogue entry
\subsection{Definitions and mathematical statement}
For a finite group \(G\), write \(|g|\) for the order of an element \(g\in G\), and use the convention \([x,y]=x^{-1}y^{-1}xy\). Define the set of \(\star\)-commutators by
\[
S(G)=\bigl\{[x,y]:x,y\in G,\ |x|=p^r,\ |y|=q^s,\ p,q\text{ distinct primes},\ r,s\in\Z_{\geq 1}\bigr\}.
\]
Thus the two elements defining a \(\star\)-commutator have coprime prime-power orders; the resulting commutator itself need not have prime-power order and may be the identity. A group \(G\) is soluble if its derived series, defined by
\[
G^{(0)}=G,\qquad G^{(n+1)}=\langle [u,v]:u,v\in G^{(n)}\rangle,
\]
satisfies \(G^{(m)}=\{1\}\) for some integer \(m\geq 0\).

\textbf{Question (V. S. Monakhov, Kourovka Notebook Problem 19.56).} Is it true that every finite group \(G\) satisfying
\[
\forall a,b\in S(G),\qquad
\gcd(|a|,|b|)=1\ \Longrightarrow\ |ab|\geq |a|\,|b|
\]
is soluble?

The supplied source does not specify a commutator convention; the convention above makes the notation explicit. Prime-power orders are interpreted as positive powers of primes. If order \(1\) is also admitted under that terminology, the only additional possible \(\star\)-commutator value is the identity, and any inequality involving it is an equality. The supplied registry labels the question open and separately records an unverified affirmative AI solution claim; that claim is not treated here as an established solution.
\subsection{Short English statement}
Must a finite group be soluble if the order of the product of any two \(\star\)-commutators of coprime orders is at least the product of their orders? Here a \(\star\)-commutator comes from two elements of coprime prime-power orders. The supplied affirmative solution claim remains explicitly unverified.
\subsection{Sources}
\begin{itemize}
\item Local UnsolvedMath registry, permanent ID \(701604\), exact source code \texttt{catalogue entry}, title \emph{Kourovka Notebook Problem 19.56}; supplied complete source record, upstream revision \texttt{d371ee917ac9647fbb06cd6fc529deec6f3dda82}.
\item V. S. Monakhov, Problem 19.56, \emph{The Kourovka Notebook (Unsolved Problems in Group Theory)}, 19th issue (2018); also attributed in the supplied record to the 21st issue edition (2026), \texttt{arXiv:1401.0300v48}. The supplied collection description names E. I. Khukhro and V. D. Mazurov as editors.
\item The supplied background and status review report an explicitly unverified affirmative AI solution claim attributed to M. Kinyon and J. Urban, \emph{Preprint} (2026): \url{https://kourovkanotebookorg.wordpress.com/wp-content/uploads/2026/09/kourovka-mathematics.pdf}. The supplied review says the claim is marked with \(\#\) in the October 2026 version of the 21st issue.
\item Supplied status review dated \texttt{2026-10-16}: canonical status \texttt{open}, evidence classification \texttt{unrefereed\_claim}. This is supplied dated curation, not an independently verified literature assessment; neither the cited notebook version nor the claimed solution has been independently checked here.
\end{itemize}
\end{document}
